\chapter*{Introducción}
\addcontentsline{toc}{chapter}{Introducción}

Este manual explica cómo utilizar la Plataforma Web Robot Explota Globos desde la perspectiva de sus usuarios. Su finalidad es orientar la consulta pública del torneo, el registro de participantes, la operación de competencias, la gestión de comunicaciones y las tareas administrativas autorizadas.

El documento evita explicaciones de programación. Cuando se menciona una palabra propia de la interfaz, como \emph{token}, \emph{panel interno} o \emph{Django Admin}, se explica por su efecto práctico para el usuario. Para detalles de instalación, código fuente, configuración del servidor o mantenimiento especializado, debe consultarse el Manual Técnico.

\section*{Alcance del manual}

El manual cubre las operaciones visibles de la plataforma:
\begin{itemize}
    \item consulta del sitio público, reglas y patrocinadores;
    \item selección de categoría y registro de robots;
    \item confirmación de asistencia mediante enlace enviado por la organización;
    \item acceso al panel interno para usuarios autorizados;
    \item operación del torneo, comunicaciones y administración, en capítulos posteriores.
\end{itemize}

\section*{Sitio público y áreas administrativas}

El sitio público puede consultarse sin iniciar sesión. Allí se encuentran la página de inicio, las reglas, los patrocinadores y el acceso al registro. Las áreas administrativas requieren un usuario autorizado y permiten ejecutar acciones que modifican información del torneo. Por eso el manual separa claramente las instrucciones para visitantes y participantes de las instrucciones para organizadores, operadores y superusuarios.

\begin{advertencia}
Las operaciones que envían datos personales, confirman asistencia o modifican resultados del torneo deben revisarse antes de aceptarse. El manual marca estas acciones como advertencias u operaciones sensibles cuando corresponde.
\end{advertencia}

\section*{Datos de ejemplo}

Las capturas incluidas en este documento utilizan datos ficticios. Los nombres, instituciones y correos visibles sirven únicamente para explicar la interfaz. En operación real, cada usuario debe ingresar información verificada y autorizada por la organización.

\section*{Pendientes institucionales}

La versión publicada mantiene como pendientes institucionales la URL definitiva de producción, el contacto oficial de soporte, el código documental institucional, firmas, aprobaciones, políticas definitivas de privacidad y respaldo, y la configuración real de comunicaciones. Estos puntos no son errores del manual; requieren confirmación formal de la organización antes de reemplazar los marcadores indicados.
